% sec-06-limitations.tex - Limitations and Future Work (full)
\section{Limitations and Future Work}

\subsection{Limitations}

The approach does not rely
on live internet search for large projects, because specifying single external URLs
is unreliable at scale; instead it depends on a curated repository of appropriately
sized inputs \cite{adoptionguide}. Because the LLM does not always have access to a
PDF compiler, some senior-author formatting (table column widths, figure-caption
spacing, and occasional page breaks) must still be applied by hand, and the paper
explicitly performs that pass in the final stage. Total input is kept under roughly
5 MB and individual files under per-file token limits, which bounds how much source
material a single build can absorb. High-dimensional or approximate inputs can make
outputs less deterministic, and oversized or overly complex inputs can surface server
errors that signal reduced quality. Most fundamentally, the method changes none of
the clinical realities: it does not alter tumor biology, shorten dose-limiting
toxicity windows or survival follow-up, or compress the fixed regulatory review
clocks, and every generated document must still pass human medical, statistical, and
regulatory review before use. The build is an independent research artifact, not an
active IND or IDE and not regulatory advice.

\subsection{Future Work}

Future work extends the method along three axes. First, coverage: applying the same
single-prompt pipeline to the Phase 2-to-3 briefing package and the Phase 3 protocol
\cite{phase2protocol}, and to the cohort-review and clinical-hold response packages
under their real clocks. Second, integration: connecting the build to live electronic
data capture and electronic document management feeds so that tables are generated
directly from validated data, and adding automated verification of generated
documents in the spirit of the verification-before-generation standard
\cite{hr9510v5}. Third, breadth: extending beyond PDAC to glioblastoma and lung
adenocarcinoma \cite{gbmsynergy, lungadeno}, and connecting to the national platform
and site-documentation pipelines so that a site can produce its full document set on
demand \cite{natplatform}. Each extension preserves the same discipline: grounded
figures, name-matched files, real-time monitorability, human in the loop.
